% Part: first-order-logic
% Chapter: introduction
% Section: substitution

\documentclass[../../../include/open-logic-section]{subfiles}

\begin{document}

\olfileid{fol}{int}{sub}

\olsection{પ્રતિસ્થાપન}

વાતો પરસ્પર કેવી રીતે જોડાય છે અને વાક્યરચના તથા અર્થવિચારનો વિકાસ
આગળની વધુ પ્રગત તપાસનો પાયો કેવી રીતે નાખે છે તે સમજાવવા એક દાખલો
ચર્ચીશું. આપણા !!{derivation} તંત્રો આપણને
$\lforall[\Obj v_0][\Atom{\Obj P}{\Obj v_0}]$માંથી
$\Atom{\Obj P}{\Obj a}$ !!{derive} કરવા દેવાં જોઈએ. કદાચ આને
અનુમાનના નિયમ તરીકે પણ જણાવવું હોય. પરંતુ એ કરવા નિયમને સૌથી સામાન્ય
શબ્દોમાં જણાવતા આવડવું જોઈએ: માત્ર $\Obj P$, $\Obj a$ અને $\Obj v_0$
માટે નહિ, પરંતુ કોઈ પણ !!{formula}~$!A$, પદ~$t$ અને !!{variable}~$x$
માટે. (યાદ રાખો કે !!{constant}s પદો છે, પરંતુ !!{constant}s અને
!!{function}sમાંથી બનેલાં વધુ જટિલ પદો પણ આપણે લઈશું.) તેથી આપણે
આવું કંઈક કહી શકવા માગીએ છીએ: ``તમે જ્યારે $\lforall[x][!A(x)]$
!!{derive}d કર્યું હોય, ત્યારે $\lforall[x]$ને દૂર કરીને અને~$x$ને
~$t$થી બદલીને મળતું પરિણામ~$!A(t)$ તારવવું વાજબી છે.'' પરંતુ
``$x$ને~$t$થી બદલો'' એટલે ચોક્કસ શું? $!A(x)$ અને~$!A(t)$ વચ્ચેનો
સંબંધ શું છે? શું આ હંમેશાં કામ કરે છે?
\sourcecorrection{OLFINT-006}{મૂળની \texttt{substitution.tex},
પંક્તિઓ~16--17માં આણ્વિક સૂત્રના બીજા દલીલસ્થાનની આસપાસ કૌંસ ખોટી
જગ્યાએ છે. અહીં અભિપ્રેત સૂત્ર
$\lforall[\Obj v_0][\Atom{\Obj P}{\Obj v_0}]$ પુનઃસ્થાપિત કર્યું છે.}

આને ચોક્કસ બનાવવા આપણે \emph{પ્રતિસ્થાપન}ની ક્રિયા વ્યાખ્યાયિત કરીએ
છીએ. પ્રતિસ્થાપન ખરેખર અટપટું છે, કારણ કે~$!A$માં આવતા બધા~$x$ને
માત્ર~$t$થી બદલી શકાતા નથી અને દરેક~$t$ને કોઈ પણ~$x$ માટે
પ્રતિસ્થાપિત કરી શકાતું નથી. ફરી અનુમાનાત્મક વ્યાખ્યાઓ વડે આ મુદ્દા
સંભાળીશું. પરંતુ આ થઈ જાય પછી ``$\lforall[x][!A(x)]$માંથી $!A(t)$
તારવો'' એ અનુમાનનિયમને ચોક્કસ રીતે કહી શકાય છે. વધુમાં, આપણે બતાવી
શકીશું કે આ એ અર્થમાં સારો અનુમાનનિયમ છે કે $\lforall[x][!A(x)]$,
~$!A(t)$ને ફલિત કરે છે. પરંતુ એ સાબિત કરવા ફરી એવી વાત સાબિત કરવી
પડશે જેને પહેલી નજરે જોઈને તમને ``આ શા માટે કરીએ છીએ?'' એવો પ્રશ્ન
થાય. $\lforall[x][!A(x)]$,~$!A(t)$ને ફલિત કરે છે એ હકીકતનો આધાર આ
વાત પર છે કે $\Sat{M}{!A(t)}$ છે કે નહિ તેનો આધાર માત્ર પદ~$t$ના
મૂલ્ય પર છે: એટલે કે, જો~$t$થી નિર્દેશાતા~$\Domain{M}$ના
!!{element}ને~$m$ કહીએ, તો $\Sat{M}{!A(t)}[s]$ ત્યારે અને માત્ર ત્યારે
જ્યારે $\Sat{M}{!A(x)}[\Subst{s}{m}{x}]$. $t$માં !!{variable}s આવે
તોપણ આ સત્ય છે, પરંતુ પરિણામને ચોક્કસ કેવી રીતે જણાવવું તેમાં સાવધ
રહેવું પડશે.

\end{document}
